\section{Билет 60. Тепловая машина. Определение и принципиальное устройство. КПД тепловой машины. Идеальная тепловая машина Карно. Цикл Карно. КПД цикла Карно (без вывода). Теорема Карно.}

\begin{definition}
\textbf{Коэффициентом полезного действия (КПД)} тепловой машины называется отношение полезной работы, совершаемой за цикл, к количеству теплоты, полученному от нагревателя:
\begin{equation}
    \eta = \frac{A}{Q_1} = \frac{Q_1 - |Q_2|}{Q_1} = 1 - \frac{|Q_2|}{Q_1}. \label{eq:efficiency_60}
\end{equation}
КПД всегда меньше единицы ($\eta < 1$), так как часть энергии неизбежно отдаётся холодильнику.
\end{definition}

\begin{definition}
\textbf{Идеальной тепловой машиной} называется машина, работающая по \textbf{циклу Карно}. Цикл Карно состоит из четырёх квазистатических (обратимых) процессов:
\begin{enumerate}
    \item \textbf{Изотермическое расширение} при температуре нагревателя $T_1$ (получение теплоты $Q_1$).
    \item \textbf{Адиабатическое расширение} (температура падает от $T_1$ до $T_2$, теплота не обменивается).
    \item \textbf{Изотермическое сжатие} при температуре холодильника $T_2$ (отдача теплоты $|Q_2|$).
    \item \textbf{Адиабатическое сжатие} (температура повышается от $T_2$ до $T_1$, возврат в исходное состояние).
\end{enumerate}
Все процессы в цикле Карно обратимы, отсутствует диссипация энергии и теплообмен при конечной разности температур.
\end{definition}

\begin{theorem}[КПД цикла Карно]
Для идеальной тепловой машины, работающей по циклу Карно между температурами нагревателя $T_1$ и холодильника $T_2$, КПД определяется формулой:
\begin{equation}
    \eta_{\text{К}} = 1 - \frac{T_2}{T_1}. \label{eq:carnot_eff_60}
\end{equation}
КПД цикла Карно зависит \textbf{только} от абсолютных температур нагревателя и холодильника и не зависит от природы рабочего тела или конструкции машины.
\end{theorem}

\begin{theorem}[Теорема Карно]
\begin{enumerate}
    \item КПД любой тепловой машины, работающей между двумя тепловыми резервуарами с температурами $T_1$ и $T_2$, не может превышать КПД идеальной машины Карно: $\eta \le \eta_{\text{К}}$.
    \item Все обратимые машины, работающие между одними и теми же температурами $T_1$ и $T_2$, имеют одинаковый КПД, равный $\eta_{\text{К}}$, независимо от рабочего вещества.
    \item КПД необратимых машин всегда строго меньше КПД Карно.
\end{enumerate}
\end{theorem}

\begin{remark}
\textbf{Физические следствия:}
\begin{itemize}
    \item Достичь $\eta = 1$ возможно только при $T_2 = 0$ К (абсолютный ноль недостижим) или $T_1 \to \infty$ (практически невозможно).
    \item Для повышения КПД реальных двигателей стремятся максимально повысить температуру нагревателя $T_1$ и понизить температуру холодильника $T_2$.
    \item Цикл Карно является теоретическим пределом эффективности преобразования теплоты в работу и служит эталоном для оценки реальных тепловых машин.
\end{itemize}
\end{remark}
